\section{Local occurrence tables for overlap windows}
\label{sec:lcs-windows}

\subsection{A short query inside a long compressed word}

The equal-bigram control in Section~\ref{sec:lcs-transitions} exposes a
different source of unnecessary work. For
\[
 T=\mathtt{aabbccacbabc},\qquad
 X=(ab)^NcT,\quad Y=(ab)^NaT,
\]
the exact longest common substring has length $2N$, but shared alphabets
and directed pairs do not prove optimality. The fallback must examine
grammar cuts. Some of its overlap queries compare a very long child
against a short tail. The previous implementation built an occurrence
table over the entire long child even though the query used only a short
suffix of it. A profile at $N=2^{500}$ places almost all charged work in
these table constructions and their local occurrence queries. Profiling
times are diagnostic only and are excluded from the timing comparison.

The correction is to localize the table's text operand as well as the
query interval. This retains the exact algorithm and its complete
fallback; it changes the grammar presented to each occurrence table.

\subsection{Two exact restrictions and one coordinate translation}

Write $\operatorname{OL}(U,V)$ for the positive lengths $k$ for which
the suffix of $U$ of length $k$ equals the prefix of $V$ of length $k$.
Set $m=\min(|U|,|V|)$. Then
\[
 \operatorname{OL}(U,V)
 =\operatorname{OL}(U[|U|-m:|U|],V[0:m]).
\]
Every feasible overlap is at most $m$, so its two compared substrings
are unchanged by these restrictions. The implementation first replaces
the operands by these equal-length words. Cache keys retain the original
root identities, while cached values remain overlap lengths and require
no translation. If the restricted words are uniform, the existing exact
uniform-word shortcut now applies even when the original roots were mixed.

For a dyadic band $\ell\le k\le h$ with
$h=\min(2\ell-1,m)$, let $P=V[0:\ell]$ and $Z=U[m-h:m]$ in the
restricted coordinates. A candidate overlap of length $k$ has its seed
starting at $m-k$ in $U$, or at $h-k$ in $Z$. Its entire seed lies in $Z$:
\[
 0\le h-k\le h-\ell,\qquad (h-k)+\ell\le h.
\]
Consequently the occurrences of $P$ in $Z$ are in bijection with precisely
the seed occurrences needed for this band. No occurrence from outside
the suffix window can contribute an overlap in the band.

The matcher therefore builds its arithmetic-progression table for $(P,Z)$
and queries the interval $[0,h]$. For every returned progression of starts
$(p,d,t)$ it uses
\[
 (p+m-h,d,t)
\]
in the existing periodic-extension test. Only the first start changes;
step and multiplicity do not. The extension test still operates on the
equal-length restricted operands and returns lengths $m-p'$. Thus both
singleton verification and periodic-run departure comparisons have exactly
the same coordinate meaning as before. The dyadic bands partition positive
overlap lengths, so the resulting progressions remain disjoint and complete.

\subsection{Cost and limits of the improvement}

Occurrence-table construction now depends on the reachable grammars of
the prefix seed and suffix window, rather than on all rules reachable
from the original text. Slices reuse complete interior subtrees and create
only boundary nodes. With grammar height $H$, the two initial restrictions
and $O(\log m)$ dyadic windows create at most $O(H\log m)$ such nodes
before accounting for reuse. This remains a polynomial local operation
on a binary straight-line program. The usual word-operation statement
must still be supplemented by polynomial costs for encoded integers.

There is no universal claim that the new reachable grammar is smaller:
boundary nodes and repeated slices can add work, and the largest band
can contain the whole text. The full-overlap control below measures that
case. All slices, tables, extension checks and cached progressions retain
the shared work and memory limits. A resource exception is still an
inconclusive search, never a negative knot verdict.

On the specified equal-bigram binary-power family, the difficult overlap
queries concern only constant-length tails. The new tables therefore
have constant size even as the long prefix grows. Root-bound preparation
and grammar navigation still depend on $\log N$. This family-specific
improvement does not prove a bound on the number of relator moves,
intermediate grammar size, or completeness of the group search. In
particular it supplies neither a general subexponential recognizer nor
a quasi-polynomial recognition theorem.

\subsection{Validation and controlled measurements}

The existing literal oracles cover 3844 binary overlap pairs, random
signed parses, exact substring witnesses, periodic critical alignments,
and complete cyclic overlap searches. Three additional tests check
asymmetric periodic operands and empty boundaries, caller-root cache
reuse, uniform overlaps surrounded by exponentially long irrelevant
context, and the exact equal-bigram answer at $N=2^{500}$ under a
20000-work allowance with expansion forbidden. The complete suite passes
690 tests with Regina 7.4.1 in 153.084 seconds.

\begin{center}\small
\begin{tabular}{@{}lrrr@{}}
\toprule
Scope & Before & A/A & Windows \\
\midrule
Small 14, summed medians & 215.475 & 215.870 & 215.012 \\
Gordian, full query & 1906.227 & 1982.828 & 1917.508 \\
\bottomrule
\end{tabular}
\end{center}
Median milliseconds; all full knot queries complete.

\begin{center}\small
\begin{tabular}{@{}llrrr@{}}
\toprule
Family & $h$ & Before & A/A & Windows \\
\midrule
Prefix & 8 & 0.118 & 0.110 & 0.110 \\
Prefix & 32 & 0.265 & 0.274 & 0.274 \\
Prefix & 500 & 3.376 & 3.385 & 3.376 \\
Interior & 8 & 0.161 & 0.159 & 0.159 \\
Interior & 32 & 0.405 & 0.397 & 0.405 \\
Interior & 500 & 4.944 & 4.946 & 4.955 \\
Equal pairs & 8 & 7.249 & 7.053 & 2.174 \\
Equal pairs & 32 & 70.293 & 70.937 & 2.487 \\
Equal pairs & 500 & limit & limit & 7.893 \\
Full overlap & 8 & 33.549 & 33.367 & 29.741 \\
Full overlap & 32 & limit & limit & limit \\
Full overlap & 500 & limit & limit & limit \\
\bottomrule
\end{tabular}
\end{center}
Median milliseconds including construction; limits are censored. Full overlap enumerates all lengths in compressed progressions, not all occurrences.

\begin{center}\small
\begin{tabular}{@{}llrrr@{}}
\toprule
Family & $h$ & Before & A/A & Windows \\
\midrule
Partial & 8 & 0.319 & 0.322 & 0.297 \\
Partial & 32 & 0.650 & 0.637 & 0.630 \\
Partial & 500 & 6.775 & 6.915 & 6.936 \\
Equal-pair cyclic & 8 & 0.514 & 0.496 & 0.501 \\
Equal-pair cyclic & 32 & 0.805 & 0.808 & 0.804 \\
Equal-pair cyclic & 500 & 6.664 & 6.586 & 6.629 \\
Quotient & 8 & 0.333 & 0.339 & 0.333 \\
Quotient & 32 & 0.983 & 0.953 & 0.942 \\
Quotient & 500 & 13.486 & 13.291 & 13.532 \\
\bottomrule
\end{tabular}
\end{center}
Median milliseconds including construction; limits are censored. Local operations do not produce knot verdicts.


For the equal-pair family, medians at $h=8$ fall from 7.249 milliseconds
(7.053 A/A) to 2.174; at $h=32$ they fall from 70.293 (70.937 A/A) to
2.487. At $h=500$ both archived controls exhaust the work allowance,
whereas localization finishes in 7.893 milliseconds with 1034 arena
nodes and 12609 charged work units. All three localized sizes use eight
window queries and 162 cumulative table cells. These are actual completed
exact LCS queries, not merely longer partial searches.

The full-overlap control compares $(ab)^N$ with itself and asks for all
positive overlap lengths, represented by arithmetic progressions. At
$h=8$, localization takes 29.741 milliseconds versus 33.549 before and
33.367 A/A, with work falling from 42506 to 37641 and nodes rising from
103 to 108. At $h=32$ and $h=500$, every arm still exhausts the allowance.
The localized $h=500$ run spends 2.639 seconds before exhaustion versus
1.177 before, and creates 2120 nodes versus 503. This is a material limit
of localization: it reaches expensive equality/periodic-extension work
that the archived table construction never reaches at the same charged
budget. The changed mix of operations has a higher wall-clock cost per
charged unit. Neither run completes, so their times do not measure a
completed-query slowdown or speedup. Future work must address that
remaining bottleneck rather than extrapolate from short-tail queries.
Section~\ref{sec:lcs-periodic} subsequently supplies a separate certified
short-period algorithm for this control; the general fallback remains needed.

All full knot queries complete with identical certificate digests and
zero calls to the changed fallback. The small-fourteen sum of medians
is 215.475 milliseconds before, 215.870 A/A and 215.012 after; Gordian
medians are 1.906, 1.983 and 1.918 seconds. The full relator operations
also have identical traces between arms. At $h=500$, the equal-pair
cyclic medians are 6.664, 6.586 and 6.629 milliseconds; successful
quotient-search medians are 13.486, 13.291 and 13.532. These controls
show no corresponding knot-corpus or relator-pipeline speedup. Timing
variation is visible in A/A and should not be promoted to an algorithmic
effect.

The comparison uses the LCS class archived at \code{34f2689e9} inside
the common current word arena, donor helpers, recognizer and certificate
replayer. A/A repeats that archived class. Seed 3012 shuffles arm order
in each round; each case has one excluded warmup and five measured rounds.
The audit includes 225 full knot queries, 180 string queries and 135
local relator operations. Full query timings include fresh PD construction
and independent replay; certificate digests are compared outside the
timer. String and relator timings include grammar construction. Local
queries allow two million work units and 100000 nodes/table cells.
Whole recognition allows 18 seconds, with 15 seconds and twenty million
work units in its group stage. Censored runs are never used as completed
runtime denominators.

The cyclic equal-bigram operation is deliberately a separate control.
Its doubled words have a common substring $T(ab)^N$ of length $2N+12$,
so the existing cyclic algorithm already finds a useful checked move
quickly. Applying it reduces the target to $c^{-1}a$, with total relator
length $2N+15$. Faster exact linear queries should not be reported as a
speedup for this already fast cyclic operation. All completed relator
traces are checked for agreement between arms; these supplied abstract
presentations are not knot-recognition verdicts.

Raw samples are in \path{results/lcs_windows_20261008.json}, with the
profile, test and build logs under \path{synthesis/data/lcs-window-*}.
The synchronization brought in fifteen research archives. Static relevance
triage inspected their overview material and scanned textual members,
with archive hashes, member names and matches retained in
\path{data/lcs-window-incoming-scan.json}. The new material concerns
recurrence avoidance, special functions, affine partitions and probabilistic
asymptotics. The affine-partition article uses perfect hash families to
coordinate rainbow zero sets, but its parameter regime is fixed field
size and growing numbers of coordinate arms; its scope statement explicitly
excludes a faster unknot recognizer. Sparse-graph moment and tree-broadcast
results concern probabilistic limits rather than the deterministic planar
diagrams or compressed group words in this implementation. No reduction
connecting these results to a useful exact recognition step was identified,
so none was incorporated as a complexity premise. This was a relevance
review, not an independent validation of those manuscripts' theorems.
The suggested \code{openai/math} head remained \code{fd4aeeb2ee}.
